(1) 一摩尔满足 $pV=RT$，因此 $T_f/T_i=(V_i/V_f)^{n-1}$。对外功 $W=\int_{V_i}^{V_f}p\,dV=K(V_f^{1-n}-V_i^{1-n})/(1-n)=(p_iV_i-p_fV_f)/(n-1)$。该式只适用于 $n\ne1$，其中 $K$ 的量纲随指数 $n$ 改变。
(2) 约定系统吸热 $Q>0$、对外做功 $W>0$。第一定律 $Q=\Delta U+W=C_V(T_f-T_i)+R(T_f-T_i)/(1-n)$，即 $Q=[C_V+R/(1-n)]\Delta T$。即使 $\Delta T$ 为正，热量的符号仍取决于路径指数，因此不可直接把温升等同吸热。
(3) 由 Mayer 关系 $C_P-C_V=R$ 得 $\gamma=1+R/C_V$。在 $n=\gamma$ 时 $C_V+R/(1-\gamma)=0$，故 $Q=0$；理想气体固定热容下等熵准静态路径正满足 $pV^\gamma=\mathrm{const}$。在 $n=0$ 时 $p$ 恒定，$W=p(V_f-V_i)=R\Delta T$、$Q=C_P\Delta T$。
(4) 令 $\epsilon=1-n$，$pV^n=K$ 使 $W=K(V_f^\epsilon-V_i^\epsilon)/\epsilon$；利用 $V^\epsilon=1+\epsilon\ln V+O(\epsilon^2)$，得到 $W\to RT\ln(V_f/V_i)$。膨胀 $V_f>V_i$ 时功为正，压缩时为负，等温下 $\Delta U=0$ 而 $Q=W$。
